% Student paper — English; IMF/diagram clips from vault assets. Compiler: XeLaTeX.
\documentclass[10pt,a4paper]{article}

\usepackage[margin=16mm,headheight=13pt,headsep=5pt,footskip=9mm]{geometry}
\usepackage{fontspec}
\usepackage[version=4]{mhchem}
\usepackage{chemfig}
\usepackage{graphicx}
\usepackage{xcolor}
\usepackage{booktabs}
\usepackage{array}
\usepackage{enumitem}
\usepackage{needspace}
\usepackage{fancyhdr}
\usepackage{amsmath,amssymb}
\usepackage{pgffor}

\raggedbottom
\setlist[enumerate]{leftmargin=1.5em,itemsep=0.12em,topsep=0.15em,parsep=0pt}
\setlist[itemize]{leftmargin=1.3em,itemsep=0.08em}

\pagestyle{fancy}
\fancyhf{}
\fancyhead[L]{\small \homeworktitle}
\fancyhead[R]{\small \homeworkmarks}
\fancyfoot[C]{\thepage}
\renewcommand{\headrulewidth}{0.3pt}

\newcommand{\homeworktitle}{Chemistry homework}
\newcommand{\homeworkmarks}{}
\newcommand{\sethomework}[2]{%
  \renewcommand{\homeworktitle}{#1}%
  \renewcommand{\homeworkmarks}{#2}%
}

\newcommand{\answerlines}[1]{%
  \par\vspace{0.08em}%
  \foreach \n in {1,...,#1}{\noindent\rule{\linewidth}{0.2pt}\\[0.62em]}%
}

\graphicspath{{images/}}
\DeclareGraphicsExtensions{.png,.jpg,.jpeg,.pdf}

\setlength{\parindent}{0pt}
\setlength{\parskip}{0.10em}
\setlength{\abovedisplayskip}{0.25em}
\setlength{\belowdisplayskip}{0.25em}

\newcommand{\qhead}[2]{%
  \par\vspace{0.28em}%
  \Needspace{4.2em}%
  \noindent{\small\textbf{#1}}\hfill{\small[#2]}\par\vspace{0.08em}%
}

\newcommand{\clipq}[3]{%
  \par
  \sbox0{\includegraphics[width=0.88\linewidth,keepaspectratio]{#3}}%
  \Needspace{\dimexpr\ht0+1.25em+2pt\relax}%
  \noindent{\small\textbf{#1}}\hfill{\small[#2]}\par\vspace{0.06em}%
  \noindent\usebox0\par
}

\newcommand{\sourcefoot}[1]{%
  {\footnotesize\color{gray}#1}\par
}

\begin{document}
\sethomework{AS Chemistry homework · 3.6 Intermolecular forces}{18 marks}

\begin{center}
{\large\bfseries AS Chemistry homework · 3.6 Intermolecular forces}\\[0.12em]
{\small CIE 9701 AS Chemistry · 18 marks · about 40 minutes}
\end{center}

{\small Topic: 3.6 Intermolecular forces, electronegativity and bond properties --- hydrogen bonding (\ce{N-H}, \ce{O-H}), van der Waals' forces (id-id, pd-pd), dipole moments, and relative strength of bonding vs intermolecular forces.
Answer in English. Section A is Paper 1 style (1 mark each). Section B is written; use syllabus terminology throughout.}

\vspace{0.15em}
{\small
\begin{tabular}{@{}llr@{}}
\toprule
\textbf{Section} & \textbf{Content} & \textbf{Marks} \\
\midrule
A & Multiple choice & 8 \\
B1 & 2023 MJ Paper 22 Question 1(b)(iii) & 1 \\
B2 & 2023 MJ Paper 22 Question 1(b)(iv) & 2 \\
B3 & 2021 MJ Paper 22 Question 2(c)(iii) & 2 \\
B4 & 2025 MJ Paper 22 Question 1(d)(i) & 3 \\
B5 & 2021 ON Paper 22 Question 1(a)(ii) & 2 \\
\midrule
 & \textbf{Total} & \textbf{18} \\
\bottomrule
\end{tabular}}

\vspace{0.25em}
{\large\bfseries Section A --- Multiple-choice questions (8 marks)}\par

\clipq{A1.}{1}{img-001.png}
\sourcefoot{2021 FM Paper 12 Question 5}

\clipq{A2.}{1}{img-002.png}
\sourcefoot{2021 MJ Paper 13 Question 5}

\clipq{A3.}{1}{img-003.png}
\sourcefoot{2024 MJ Paper 11 Question 8}

\clipq{A4.}{1}{img-004.png}
\sourcefoot{2024 MJ Paper 12 Question 8}

\clipq{A5.}{1}{img-005.png}
\sourcefoot{2021 ON Paper 12 Question 6}

\clipq{A6.}{1}{img-006.png}
\sourcefoot{2021 ON Paper 11 Question 33}

\clipq{A7.}{1}{img-007.png}
\sourcefoot{2021 MJ Paper 12 Question 3}

\clipq{A8.}{1}{img-008.png}
\sourcefoot{2025 MJ Paper 12 Question 10}

\vspace{0.35em}
{\large\bfseries Section B --- Structured questions (10 marks)}\par\vspace{0.1em}

% B1/B2 paper clips wrongly included the Paper 2 cover page + an unrelated melting-point table.
\qhead{B1.}{1}
Name the strongest intermolecular force that exists between \ce{NH3(l)} molecules.
\answerlines{1}
\sourcefoot{2023 MJ Paper 22 Question 1(b)(iii)}

\qhead{B2.}{2}
Draw a diagram to show the formation of the strongest intermolecular force between \textbf{two} molecules of \ce{NH3(l)}.

Include any relevant lone pairs of electrons and dipoles.

\vspace{0.2em}
\noindent\fbox{\parbox[c][4.4cm][c]{0.92\linewidth}{\centering\vspace{3.6cm}}}
\par\vspace{0.15em}
\sourcefoot{2023 MJ Paper 22 Question 1(b)(iv)}

% B3 clip started with leftover words from the previous part.
\qhead{B3.}{2}
Nitrogen, \ce{N2}, is a gas at room temperature and pressure. Neither \ce{CO} nor \ce{N2} is an ideal gas.

Complete the table by naming \textbf{all} the types of intermolecular forces (van der Waals') in separate samples of \ce{N2(g)} and \ce{CO(g)}.

\vspace{0.15em}
{\small
\noindent\begin{tabular}{|p{0.42\linewidth}|c|c|}
\hline
 & \ce{N2(g)} & \ce{CO(g)} \\
\hline
number of electrons per molecule & 14 & 14 \\
\hline
presence of a dipole moment & $\times$ & $\checkmark$ \\
\hline
boiling point / $^\circ$C & $-195.8$ & $-191.5$ \\
\hline
intermolecular forces (van der Waals') & \rule{0pt}{2.6em} &  \\
\hline
\end{tabular}}
\par\vspace{0.1em}
\sourcefoot{2021 MJ Paper 22 Question 2(c)(iii)}

\qhead{B4.}{3}
Table 1.1 shows melting points of some oxides.

Complete Table 1.1 by identifying the strongest force of attraction in each oxide that is broken during melting. Use the abbreviations below.

{\small
\begin{itemize}[nosep]
\item i.d.\ = instantaneous dipole--induced dipole
\item p.d.\ = permanent dipole--permanent dipole
\item H = hydrogen bond
\item C = covalent bond
\item I = ionic bond
\end{itemize}}

\vspace{0.1em}
{\small
\noindent\begin{tabular}{|c|c|p{0.42\linewidth}|}
\hline
oxide & melting point / $^\circ$C & force of attraction broken during melting \\
\hline
\ce{SO2} & $-73$ & \rule{0pt}{1.35em} \\
\hline
\ce{H2O} & 0 &  \\
\hline
\ce{SiO2} & 1610 &  \\
\hline
\ce{Na2O} & 1132 &  \\
\hline
\ce{MgO} & 2852 &  \\
\hline
\end{tabular}}
\par\vspace{0.1em}
\sourcefoot{2025 MJ Paper 22 Question 1(d)(i)}

\clipq{B5.}{2}{img-013.png}
\sourcefoot{2021 ON Paper 22 Question 1(a)(ii)}

\end{document}
